\section{Getting Started}
\label{sec:GettingStarted}

This section walks through building and running Archimesh locally. It assumes
Docker Desktop (or Docker Engine with the Compose plugin) and Java~21+.

\subsection{Prerequisites}

\begin{itemize}
\item Docker Desktop (or Docker Engine + Compose plugin)
\item Java 21+ (for the Quarkus server)
\item Maven 3.9+ (for building the server and client plugin)
\item Eclipse with Archi and the Archi modelling tool (for the client plugin)
\end{itemize}

\subsection{Configuration}

\begin{enumerate}
\item Copy \filepath{.env.example} to \filepath{.env}:
\begin{lstlisting}[language=bash,caption={Setup environment}]
cp .env.example .env
\end{lstlisting}

\item Edit \filepath{.env} if you need to change the Neo4j password or Kafka
  bootstrap address. The defaults work for local development:
\begin{lstlisting}[caption={Default .env values}]
NEO4J_USER=neo4j
NEO4J_PASSWORD=devpassword
NEO4J_URI=bolt://localhost:7687
KAFKA_BOOTSTRAP_SERVERS=localhost:9092
KAFKA_TOPIC_PREFIX=archimesh.model
KAFKA_CONSUMER_GROUP_ID=archimesh-server
ARCHIMESH_CONSISTENCY_CHECK=false
ARCHIMESH_AUTHZ_ENABLED=false
ARCHIMESH_IDENTITY_MODE=bootstrap
ARCHIMESH_LOG_LEVEL=INFO
\end{lstlisting}
\end{enumerate}

\subsection{Starting infrastructure}

Start Neo4j and Kafka with Docker Compose:

\begin{lstlisting}[language=bash,caption={Start containers}]
docker compose up -d
\end{lstlisting}

Services available after startup:

\begin{description}
\item[Neo4j Browser] \url{http://localhost:7474}
\item[Neo4j Bolt] \texttt{bolt://localhost:7687}
\item[Kafka broker] \texttt{localhost:9092}
\item[Kafka UI] \url{http://localhost:8080}
\end{description}

Apply Neo4j constraints and indexes:

\begin{lstlisting}[language=bash,caption={Apply schema}]
docker compose exec -T neo4j cypher-shell \
  -u "${NEO4J_USER:-neo4j}" \
  -p "${NEO4J_PASSWORD:-devpassword}" \
  < neo4j/schema.cypher
\end{lstlisting}

\subsection{Building and running the server}

The server is a Quarkus application in \filepath{server/}:

\begin{lstlisting}[language=bash,caption={Build and run server}]
cd server
mvn quarkus:dev
\end{lstlisting}

The server starts on port 8081. Verify:

\begin{lstlisting}[language=bash,caption={Health check}]
curl -s http://localhost:8081/admin/overview | jq .
\end{lstlisting}

\subsection{Creating a model}

Models must be explicitly registered before use:

\begin{lstlisting}[language=bash,caption={Register model}]
curl -s -X POST http://localhost:8081/admin/models/demo | jq .
\end{lstlisting}

\subsection{Building the client plugin}

The client plugin is an Eclipse RCP bundle:

\begin{lstlisting}[language=bash,caption={Build client plugin}]
cd client
mvn -q -pl org.gautelis.archimesh.plugin -am install -DskipTests
\end{lstlisting}

Install the built plugin into Archi using:

\begin{lstlisting}[language=bash,caption={Install plugin}]
./scripts/install-archimesh-plugin.sh
\end{lstlisting}

\subsection{Building the admin UI}

The admin UI is a Svelte application:

\begin{lstlisting}[language=bash,caption={Build admin UI}]
cd admin-web
npm install
npm run build
\end{lstlisting}

Or use the convenience script:

\begin{lstlisting}[language=bash,caption={Build admin UI via script}]
./scripts/build-admin-web.sh
\end{lstlisting}

The built static app is served by Quarkus at \texttt{/admin-ui/}.

\subsection{First collaborative session}

\begin{enumerate}
\item Start two Archi instances with the plugin installed.
\item In the first instance, create a new model or open an existing one.
\item Use \textbf{Tools \textgreater\ Connect Archimesh...} to connect.
  Enter the WebSocket URL (e.g. \texttt{ws://localhost:8081}) and model id
  (\texttt{demo}).
\item In the second instance, connect to the same model.
\item Make edits in either instance. Changes appear in the other within the
  broadcast latency of the Kafka fan-out path (typically under one second).
\end{enumerate}

\subsection{Validation scripts}

The project includes local validation scripts for fast feedback:

\begin{lstlisting}[language=bash,caption={Validation commands}]
# Fast hardening and parity checks (no Docker infra needed)
./scripts/validate-local.sh fast

# Full client + server gate
./scripts/validate-local.sh gate

# Full gate with local Kafka/Neo4j-backed integration tests
./scripts/validate-local.sh infra

# Full gate with websocket end-to-end coverage
./scripts/validate-local.sh ws
\end{lstlisting}

These are the supported release gate entrypoints. Run \texttt{fast} before
committing; run \texttt{gate} or \texttt{infra} before pushing.

\subsection{Stopping}

\begin{lstlisting}[language=bash,caption={Stop infrastructure}]
docker compose down
\end{lstlisting}

Add \texttt{-v} to also remove volumes (clean slate):

\begin{lstlisting}[language=bash,caption={Stop and remove volumes}]
docker compose down -v
\end{lstlisting}

\subsection{Where to go next}

\begin{itemize}
\item Section~\ref{sec:Architecture} for the system architecture and component
  layout.
\item Section~\ref{sec:ClientPlugin} for the client plugin internals.
\item Section~\ref{sec:ServerApi} for the complete API reference.
\item Section~\ref{sec:AdminOperations} for the admin UI and supervision
  workflows.
\item \filepath{dev-setup.md} for additional local development tips.
\end{itemize}
